% from main.tex
%%%%%%%%%%  scan idx 114  |  folio 101  %%%%%%%%%%
est
pleinement fidèle et a pour image essentielle \(\Dpx(\mathcal C)_{\mathcal C_0\text{-}\mathrm{coh}}\).

\emph{Preuve.} a) Soit \(F\in\Ob\Dpx(\mathcal C)_{\mathcal C_0\text{-}\mathrm{coh}}\), et supposons construit un \(i\)-quasi-isomorphisme (de \(\Cpx(\mathcal C)\)) \(u_i:E_i\longrightarrow F\) avec \(E_i\) strictement pseudo-cohérent (c'est possible pour \(i\) assez grand avec \(E_i=0\)). Soit \(G_i\) le cône de \(u_i\). Il résulte de (1.4.2) et (2.1.1) que \(\Ho^{i-1}(G_i)\) est de \(\mathcal C_0\)-type fini. Donc, en vertu de (1.4.4), il existe un triangle commutatif de \(\Cpx(\mathcal C)\)
\[
\begin{tikzcd}
E_i \arrow[r,"v_{i-1}"] \arrow[d,"u_i"'] & E_{i-1} \arrow[dl,"u_{i-1}"]\\
F &
\end{tikzcd}
\]
où \(E_{i-1}\) est strictement pseudo-cohérent, \(u_{i-1}\) un \((i-1)\)-quasi-isomorphisme, et \(v_{i-1}\) un monomorphisme induisant l'identité en degré \(\geq i\). On en déduit par passage à la limite (cf. preuve de (1.4.5)) un quasi-isomorphisme \(u:E\to F\) avec \(E\) strictement pseudo-cohérent.

b) Appliquer le critère ([V] chap I § 2 4.2 b)(i)) et a).

Noter que la preuve de a) donne le résultat plus précis suivant, qui généralise (1.4.5) :

\textbf{Proposition 2.7.1.} Sous les hypothèses de (2.7), soit \(u:E\to F\) un \(n\)-quasi-isomorphisme de \(\Cpx(\mathcal C)\), avec \(F\) pseudo-cohérent et \(E\) strictement pseudo-cohérent. Il existe alors un triangle commutatif de \(\Cpx(\mathcal C)\)
\[
\begin{tikzcd}
E \arrow[r,"v"] \arrow[d,"u"'] & E' \arrow[dl,"u'"]\\
F &
\end{tikzcd}
\]
où \(E'\) est strictement pseudo-cohérent, \(u'\) un quasi-isomorphisme, et \(v\) un
